\documentclass[12pt]{beamer}
% \usetheme{umbc4}
\usetheme{moloch}

\usepackage{amsmath,amscd}
\usepackage{fontspec}
\usepackage{unicode-math}
\usepackage{pifont}
\usepackage{xcolor}
\usepackage[dvipsnames]{xcolor}
\usepackage{listings}
\usepackage{minted}
\usepackage{tikz}
\usepackage{twemojis}
\usepackage{nicematrix}

\usetikzlibrary{fit}

\setmonofont{DejaVu Sans Mono}

% Set the title bar to black background with white text
\setbeamercolor{frametitle}{bg=black, fg=white}

% Ensure other structural elements (bullets, etc.) are visible
\setbeamercolor{structure}{fg=white}

\setmainfont{Libertinus Serif}
\setsansfont{Libertinus Sans}
\setmathfont{STIX Two Math}

% Define colors for syntax highlighting
\definecolor{codegreen}{rgb}{0,0.6,0}
\definecolor{codegray}{rgb}{0.5,0.5,0.5}
\definecolor{codecoffee}{rgb}{0.75,0.6,0.3}
\definecolor{question}{rgb}{1.0,0.75,0.1}
\definecolor{lightblue}{rgb}{0.55,0.75,0.9}
\definecolor{yellow}{rgb}{1.0, 1.0, 0.0}
\definecolor{codecomment}{rgb}{0.97, 0.74, 0.4}
\definecolor{codepy}{rgb}{0.9, 0.85, 0.85} %{#F8BC65}
\definecolor{funcColor}{rgb}{0.47, 0.6, 0.8}
\definecolor{midgray}{rgb}{0.4, 0.4, 0.3}
\definecolor{indexwhite}{rgb}{0.9, 0.9, 0.9}

\setbeamercolor{background canvas}{bg=black}
\setbeamercolor{normal text}{fg=white}


\newcommand{\FrameTitle}[2]{\frametitle{#1 \hfill \small\textnormal{\textcolor{midgray}{#2}}}}
\newcommand{\snl}{\vspace*{8pt}\\}
\newcommand{\nl}{\vspace*{1em}\\}
\newcommand{\vsp}{\vspace*{1em}}

\newcommand{\Matrix}[1]{\textbf{#1}}

\begin{document}
\title{\center{Applied Linear Algebra \\\vsp Vector Spaces \\}}

\author{\center{Devi Prasad M \\ Manipal School of Information Sciences \\ Manipal \\}}
\date{\center{Monday, 3 August 2026}}
\maketitle


\begin{frame}\FrameTitle{Scalars}{Unit 1}
    \begin{itemize}
        \item A scalar is a single value, which can be an integer, a real number, or a complex number.
        \item Scalars are the simplest mathematical objects.
    \end{itemize}
    \vsp
    {\color{question}{\Large\textbf{Quiz:}}} \snl {\color{white}{\large{Name a few good examples of scalars.}}}
\end{frame}

\begin{frame}\FrameTitle{Vectors}{Unit 1}
    \begin{itemize}
        \item A vector is a more structured mathematical object.
        \item A vector is an ordered sequence of scalars.\snl
        \item[] \textbf{v} $ = (a_1, a_2, \ldots, a_n) \in \mathbb{R}^n$
        \item[] \hspace*{1em} $where \ \mathbb{R}^n$ is the set of all n-dimensional real vectors. \snl
        \item[]$ v = [0.35, 0.71, 0.23, 0.8] \in \mathbb{R}^4$ \snl
        \item[]$ v = [0.35] \in \mathbb{R}^{\textbf{\color{yellow}{?}}}$ \snl
        \item[]$ v = [0.35, 0.71] \in \mathbb{R}^{\textbf{\color{yellow}{?}}}$ \snl
        \item[]$ v = [] \in \mathbb{R}^{\textbf{\color{yellow}{?}}}$ \snl
    \end{itemize}
\end{frame}

\begin{frame}\FrameTitle{Vectors}{Unit 1}
    Let \Matrix{A} be a $m \times n$ matrix.\snl
    The $n$ columns of \Matrix{A} are vectors in $\mathbb{R}^m$.\snl

    \[
    \begin{pNiceMatrix}[
        margin=3pt,
        code-after = {
            \tikz[remember picture, overlay]
            \node [
                draw=yellow,          % Soft blue border line
                % fill=blue!10,          % Very light blue interior fill
                opacity=0.9,           % High transparency for the shape
                text opacity=1,        % Keeps numbers dark and 100% readable
                thick,                  % outline border style
                rounded corners=5pt,   % Smooths the rectangle corners perfectly
                inner xsep=4pt,        % Padding to prevent numbers from touching the border
                inner ysep=6pt,        % Vertical padding at the top and bottom
                fit = (1-2) (5-2)      % Fits from Row 1, Col 2 down to Row 5, Col 2
            ] {} ;
        }
    ]
    .1 & .2 & .3 & 0.6\\
    .4 & .5 & .6 & 0.12\\
    .7 & .8 & .9 & 0.18\\
    .11 & .12 & .13 & 0.26\\
    .21 & .22 & .23 & 0.33
    \end{pNiceMatrix}
    \]
\end{frame}



\begin{frame}\FrameTitle{Vectors}{Unit 1}
    A vector can be represented as a row or column of scalars.\\
    \begin{columns}[T]
        \begin{column}{0.48\textwidth}
            \centering{\color{lightblue}{Row vector}}
            \[
            \begin{bmatrix}
                0.35 & 0.71 & 0.23 & 0.8
            \end{bmatrix}
            \]
            \[
            \begin{pmatrix}
                0.35 & 0.71 & 0.23 & 0.8
            \end{pmatrix}
            \]
        \end{column}
        \vline{\color{darkgray}\vrule width 0.2pt}
        \begin{column}{0.48\textwidth}
            \centering{\color{lightblue}{Column vector}}
            \[
            \begin{bmatrix}
                0.35 \\ 0.71 \\ 0.23 \\ 0.8
            \end{bmatrix}
            \]
            \[
            \begin{pmatrix}
                0.35 \\ 0.71 \\ 0.23 \\ 0.8
            \end{pmatrix}
            \]
        \end{column}
    \end{columns}
\end{frame}

\begin{frame}\FrameTitle{The Zero-Dimensional Vector Space $\mathbb{R}^0$}{Unit 1}
    % 1. Definition Block
    \begin{definition}[The Trivial Space]
        The space $\mathbb{R}^0$ is the unique zero-dimensional real vector space.
        It contains exactly one element, the empty vector $v$:
        \[
        v \in \mathbb{R}^0 \quad \text{where} \quad v = [ ]
        \]
    \end{definition}

\end{frame}

\begin{frame}\FrameTitle{Vectors}{Unit 1}
    {\color{question}{\Large\textbf{Quiz:}}} \snl {\color{white}{\large{Name a few good examples of vectors.}}}
\end{frame}

\begin{frame}\FrameTitle{\color{question}{Math vs. Code}}{\twemoji{light bulb}}
    --- Textbooks use \emph{column vectors} for single samples.\snl
    --- {\color{midgray}{PyTorch uses \emph{row vectors} to handle batches.}}\snl

    \begin{itemize}
        \item Single sample $\mathbf{x}$ is a column vector.
        \item Forward pass formula:
            \[
            \mathbf{y} = \mathbf{W}\mathbf{x} + \mathbf{b}
            \]
            \item Dimensions:
            \[
            \begin{matrix}
            \mathbf{y} & = & \mathbf{W} & \mathbf{x} & + & \mathbf{b} \\
            {\color{yellow}{\scriptstyle(c \times 1)}} & & {\color{yellow}{\scriptstyle(c \times d)}} & {\color{yellow}{\scriptstyle(d \times 1)}} & & {\color{yellow}{\scriptstyle(c \times 1)}}
            \end{matrix}
            \]
    \end{itemize}
\end{frame}

\begin{frame}\FrameTitle{\color{question}{Math vs. Code}}{\twemoji{light bulb}}
    --- {\color{midgray}{Textbooks use \emph{column vectors} for single samples.}}\snl
    --- PyTorch uses \emph{row vectors} to handle batches.\snl
    \begin{itemize}
        \item Batch $\mathbf{X}$ stacks samples as rows.
        \item PyTorch \emph{nn.Linear} formula:
        \[
        \mathbf{Y} = \mathbf{X}\mathbf{W}^T + \mathbf{B}
        \]
        \item Dimensions ($N$ = batch size):
        \[
        \begin{matrix}
        \mathbf{Y} & = & \mathbf{X} & \mathbf{W}^T & + & \mathbf{B} \\
        {\color{yellow}{\scriptstyle(N \times c)}} & & {\color{yellow}{\scriptstyle(N \times d)}} & {\color{yellow}{\scriptstyle(d \times c)}} & & {\color{yellow}{\scriptstyle(1 \times c)}}
        \end{matrix}
        \]
    \end{itemize}
\end{frame}

\begin{frame}\FrameTitle{\color{question}{NumPy and PyTorch}}{\twemoji{light bulb}}
    The standard convention for data analysis and machine learning is that
    \begin{itemize}
        \item[---] rows represent samples or observations: \textcolor{brown}{\textbf{axis 0}}
        \item[---] columns represent features or attributes: \textcolor{brown}{\textbf{axis 1}}
    \end{itemize}
\end{frame}

\begin{frame}\FrameTitle{Vector Spaces}{Unit 1}
A vector space $V$ is a set endowed with a rule for addition and a
rule for scalar multiplication
    \begin{align*}
        +: V \times V \rightarrow V \\
        \cdot: \mathbb{R} \times V \rightarrow V
    \end{align*}

such that the following axioms hold for all $u, v, w \in V$ and all scalars $c_{1}, c_{2} \in \mathbb{R}$:
    \begin{enumerate}
        \item $u + v = v + u$
        \item $(u + v) + w = u + (v + w)$
        \item $v + \mathbf{0} = v$
        \item $v + (-v) = \mathbf{0}$ where $-v$ is the unique additive inverse of $v$
        \item $c_1 \cdot (v + w) = c_1 \cdot v + c_1 \cdot w$
        \item $(c_1 + c_2) \cdot v = c_1 \cdot v + c_2\cdot v$
        \item $c_1 \cdot (c_2 \cdot v) = (c_1 c_2) \cdot v$
        \item $1 \cdot v = v$
    \end{enumerate}
\end{frame}

\begin{frame}\FrameTitle{Vector Spaces}{Unit 1}
    {\center{\Large\textbf{Examples of Vector Spaces}}}
\end{frame}

\begin{frame}\FrameTitle{Vector Spaces - Quiz}{Unit 1}
    \color{question}{\Large\textbf{Question:}} \snl
    {\color{white}{\large{is a $3 \times 2$ matrix of $\mathbb{R}$ a vector space?}}}
    \snl
    {\color{white}{\large{is a $3 \times 2$ matrix with non-negative entries a vector space?}}}
    \snl
    {\color{white}{\large{is a $m \times n$ matrix of $\mathbb{R}$ a vector space?}}}
\end{frame}

\begin{frame}\FrameTitle{Vector Spaces - $\mathbb{R}^1, \mathbb{R}^2, \mathbb{R}^3,\ldots,\mathbb{R}^n$}{Unit 1}
    \begin{definition}
        The space $\mathbb{R}^n$ consists of all column vectors with $n$ components,
        where each component is a real number.
    \end{definition}

    \vsp
    {\Large\textbf{\color{question}Quiz}}
    \vsp
    \begin{enumerate}
    \item {\textcolor{question}{
        Suppose addition in $\mathbb{R}^2$ adds an extra 1 to each component.
        With scalar multiplication unchanged, which rules are broken?
    }}
    \nl
    \item {\textcolor{question}{
        Show that the set of all positive real numbers, with $x + y$ and $cx$ redefined to
        equal the usual $xy$ and $x^c$, is a vector space. What is the ``zero vector''?
    }}
    \end{enumerate}
\end{frame}

\begin{frame}\FrameTitle{Subspaces}{Unit 1}
    \begin{definition}
        A subspace $W$ of a vector space $V$ is a non-empty subset of $V$
        that is itself a vector space under the same operations of addition
        and scalar multiplication defined on $V$.
    \end{definition}
    \begin{enumerate}
        \item The zero vector of $V$ is in $W$.
        \item If $u, v \in W$, then $u + v \in W$ (closed under addition).
        \item If $c$ is a scalar and $v \in W$, then $cv \in W$ (closed under scalar multiplication).
    \end{enumerate}
\end{frame}

\begin{frame}\FrameTitle{Vector subspaces}{Unit 1}
    Consider all vectors in $\mathbb{R}^2$ whose components are positive or zero. \snl
    This subset is the first quadrant of the x-y plane.\snl
    It is not a subspace of $\mathbb{R}^2$ because it is not closed under scalar multiplication.\snl
    For example, if
    $$
        v = [1, 1]
    $$
    is in the first quadrant,
    then
    $$
        -1 \cdot v = [-1, -1]
    $$
    is not in the first quadrant.
\end{frame}

\begin{frame}\FrameTitle{Subspaces}{Unit 1}
    If we include the third quadrant along with the first, scalar multiplication is closed,
    but the set is still not a subspace because it is not closed under addition. \snl

    For example, if
    $$
        u = [1, 2] \quad \text{and} \quad v = [-2, -1]
    $$
    are in the first and third quadrants, then
    $$
        u + v = [-1, 1]
    $$
    is not in the first or third quadrant.
\end{frame}

\begin{frame}\FrameTitle{Subspaces}{Unit 1}
    The smallest subspace of $\mathbb{R}^2$ that contains
    the first quadrant is the entire space $\mathbb{R}^2$.
\end{frame}

\end{document}